{\rtf1\ansi\ansicpg1252\cocoartf2639
\cocoatextscaling0\cocoaplatform0{\fonttbl\f0\fswiss\fcharset0 Helvetica;}
{\colortbl;\red255\green255\blue255;}
{\*\expandedcolortbl;;}
\paperw11900\paperh16840\margl1440\margr1440\vieww11520\viewh8400\viewkind0
\pard\tx566\tx1133\tx1700\tx2267\tx2834\tx3401\tx3968\tx4535\tx5102\tx5669\tx6236\tx6803\pardirnatural\partightenfactor0

\f0\fs24 \cf0 \\documentclass[11pt]\{article\}\
\
% ===================== PACKAGES =====================\
\\usepackage[a4paper,margin=1in]\{geometry\}\
\\usepackage\{amsmath,amssymb\}\
\\usepackage\{enumitem\}\
\\usepackage\{fancyhdr\}\
\\usepackage\{xcolor\}\
\
% ===================== HEADER & FOOTER =====================\
\\pagestyle\{fancy\}\
\\fancyhf\{\}\
\\lhead\{CSE 400: Fundamentals of Probability in Computing\}\
\\rhead\{Lecture Scribe -- Lecture 5\}\
\\cfoot\{\\thepage\}\
\
% ===================== TITLE =====================\
\\title\{\
    \\normalsize School of Engineering and Applied Science (SEAS), Ahmedabad University \\\\\
    \\vspace\{0.2cm\}\
    \\textbf\{CSE 400: Fundamentals of Probability in Computing\}\\\\\
    \\Large Lecture 5 Scribe: Bayes\'92 Theorem, Random Variables, and Probability Mass Function\
\}\
\\author\{\}\
\\date\{\}\
\\begin\{document\}\
\\maketitle\
\
\\vspace\{-2cm\}\
\\begin\{center\}\
    \\begin\{tabular\}\{ll\}\
        \\textbf\{Instructor:\} & \{Prof. Dhaval Patel\} \\\\ [1.5ex]\
        \\textbf\{Course:\} & \{CSE 400 -- Fundamentals of Probability in Computing\} \\\\ [1.5ex]\
        \\textbf\{Lecture:\} & \{Lecture 5\} \\\\ [1.5ex]\
        \\textbf\{Date:\} & \{20th Jan 2026\} \\\\ [1.5ex]\
    \\end\{tabular\}\
\\end\{center\}\
\
\\hrule\
\\vspace\{0.5cm\}\
\
% ===================== LECTURE SCRIBE =====================\
\
\\section\{Bayes\'92 Theorem\}\
\
\\subsection\{Weighted Average of Conditional Probabilities\}\
\
Let $A$ and $B$ be events. We may express\
\\[\
A = AB \\cup AB^c\
\\]\
For an outcome to be in $A$, it must either be in both $A$ and $B$ or be in $A$ but not in $B$.\
\
Since $AB$ and $AB^c$ are mutually exclusive, by Axiom 3,\
\\[\
\\Pr(A) = \\Pr(AB) + \\Pr(AB^c)\
\\]\
\
Using conditional probability,\
\\[\
\\Pr(A) = \\Pr(A \\mid B)\\Pr(B) + \\Pr(A \\mid B^c)\\,[1-\\Pr(B)]\
\\]\
\
\\textbf\{Statement.\}  \
The probability of event $A$ is a weighted average of the conditional probabilities, with weights given by the probability of the conditioning event.\
\
\\subsection\{Learning by Example\}\
\
\\textbf\{Example 3.1 (Part 1).\}\
\
An insurance company divides people into two classes:\
\\begin\{itemize\}\
    \\item Accident prone\
    \\item Not accident prone\
\\end\{itemize\}\
\
\\[\
\\Pr(\\text\{accident in 1 year\} \\mid \\text\{accident prone\}) = 0.4\
\\]\
\\[\
\\Pr(\\text\{accident in 1 year\} \\mid \\text\{not accident prone\}) = 0.2\
\\]\
\\[\
\\Pr(\\text\{accident prone\}) = 0.3\
\\]\
\
Let $A_1$ denote the event that the policyholder has an accident within one year, and let $A$ denote the event that the policyholder is accident prone.\
\
\\[\
\\Pr(A_1) = \\Pr(A_1 \\mid A)\\Pr(A) + \\Pr(A_1 \\mid A^c)\\Pr(A^c)\
\\]\
\\[\
= (0.4)(0.3) + (0.2)(0.7) = 0.26\
\\]\
\
\\textbf\{Example 3.1 (Part 2).\}\
\
Given that a policyholder has an accident within a year, find $\\Pr(A \\mid A_1)$.\
\
\\[\
\\Pr(A \\mid A_1) = \\frac\{\\Pr(A)\\Pr(A_1 \\mid A)\}\{\\Pr(A_1)\}\
= \\frac\{(0.3)(0.4)\}\{0.26\} = \\frac\{6\}\{13\}\
\\]\
\
\\section\{Law of Total Probability and Bayes Formula\}\
\
Suppose $B_1, B_2, \\dots, B_n$ are mutually exclusive events such that\
\\[\
\\bigcup_\{i=1\}^\{n\} B_i = B\
\\]\
\
Then\
\\[\
A = \\bigcup_\{i=1\}^\{n\} AB_i\
\\]\
\
Since $AB_i$ are mutually exclusive,\
\\[\
\\Pr(A) = \\sum_\{i=1\}^\{n\} \\Pr(AB_i)\
= \\sum_\{i=1\}^\{n\} \\Pr(A \\mid B_i)\\Pr(B_i)\
\\]\
\
This is the \\textbf\{Law of Total Probability (Formula 3.4)\}.\
\
\\subsection\{Bayes Formula\}\
\
Using\
\\[\
\\Pr(AB_i) = \\Pr(B_i \\mid A)\\Pr(A)\
\\]\
\
we obtain\
\\[\
\\Pr(B_i \\mid A) =\
\\frac\{\\Pr(A \\mid B_i)\\Pr(B_i)\}\
\{\\sum_\{i=1\}^\{n\} \\Pr(A \\mid B_i)\\Pr(B_i)\}\
\\]\
\
This is known as the \\textbf\{Bayes Formula (Proposition 3.1)\}.\
\
\\begin\{itemize\}\
    \\item $\\Pr(B_i)$: a priori probability\
    \\item $\\Pr(B_i \\mid A)$: posteriori probability of $B_i$ given $A$\
\\end\{itemize\}\
\
\\section\{Additional Example\}\
\
\\textbf\{Example 3.2 (Cards Problem).\}\
\
Three cards: RR, BB, RB.  \
Observed: upper side is red.\
\
\\[\
\\Pr(RB \\mid R) =\
\\frac\{(1/2)(1/3)\}\
\{(1)(1/3)+(1/2)(1/3)+(0)(1/3)\} = \\frac\{1\}\{3\}\
\\]\
\
Hence,\
\\[\
\\Pr(\\text\{other side is black\}) = \\frac\{1\}\{3\}\
\\]\
\
The lecture explains that assuming probability $1/2$ is incorrect because outcomes are not equally likely.\
\
\\section\{Random Variables\}\
\
A random variable is a real-valued function defined on the sample space:\
\\[\
X : \\Omega \\rightarrow \\mathbb\{R\}\
\\]\
\
The lecture restricts attention to discrete random variables.\
\
\\subsection\{Distribution\}\
\
For any value $a$,\
\\[\
\\Pr[X = a]\
\\]\
is the probability of the event $\\\{\\omega \\in \\Omega : X(\\omega)=a\\\}$.\
\
\\section\{Example\}\
\
Tossing 3 fair coins.  \
Let $Y$ be the number of heads.\
\
\\[\
\\Pr(Y=0)=\\frac\{1\}\{8\}, \\;\
\\Pr(Y=1)=\\frac\{3\}\{8\}, \\;\
\\Pr(Y=2)=\\frac\{3\}\{8\}, \\;\
\\Pr(Y=3)=\\frac\{1\}\{8\}\
\\]\
\
\\[\
\\sum_\{i=0\}^\{3\} \\Pr(Y=i) = 1\
\\]\
\
\\section\{Probability Mass Function\}\
\
Let $X$ be discrete with range $R_X=\\\{x_1,x_2,\\dots\\\}$.\
\
\\[\
p(x_k)=\\Pr(X=x_k)\
\\]\
\
\\[\
\\sum_k p(x_k)=1\
\\]\
\
\\section\{PMF Example\}\
\
\\[\
p(i)=c\\frac\{\\lambda^i\}\{i!\}, \\quad i=0,1,2,\\dots\
\\]\
\
\\[\
c\\sum_\{i=0\}^\{\\infty\}\\frac\{\\lambda^i\}\{i!\}=1\
\\Rightarrow c=e^\{-\\lambda\}\
\\]\
\
\\[\
\\Pr(X=0)=e^\{-\\lambda\}\
\\]\
\
\\[\
\\Pr(X>2)=1-\\Pr(X=0)-\\Pr(X=1)-\\Pr(X=2)\
\\]\
\
\\section\{Class Participation\}\
\
Students were instructed to switch to Campuswire for participation and discussion.\
\
\\bigskip\
\\hrule\
\\vspace\{0.2cm\}\
\\begin\{center\}\
    \\small \\textit\{End of Lecture 5 Scribe\}\
\\end\{center\}\
\
\\end\{document\}\
}